\section{Propositions}

% =========================================================
% TOPIC TITLE
% =========================================================

\begin{frame}{Propositions}

\begin{center}
\vspace{1cm}
{\Huge\textbf{Propositions}}

\vspace{0.6cm}

{\Large Representing and evaluating statements in computing contexts}

\end{center}

\end{frame}

% =========================================================
% QUESTION 1
% =========================================================

\begin{frame}{Propositions -- Question 1}

\begin{block}{Question}

A software system checks whether a user can access a protected
database. The following conditions are defined:

\begin{itemize}
\item $P$: The user has entered a valid password.
\item $Q$: The user has completed two-factor authentication.
\end{itemize}

The system grants access only when both conditions are satisfied.

\medskip

\begin{enumerate}
\item Express the access rule as a proposition.
\item Determine whether the user is granted access if
$P$ is true and $Q$ is false.
\end{enumerate}

\end{block}

\end{frame}

% ---------------------------------------------------------
% ANSWER 1 -- PART 1
% ---------------------------------------------------------

\begin{frame}{Propositions -- Question 1: Answer}

\begin{exampleblock}{Step 1: Represent the Rule}

The system grants access only when both conditions are satisfied.

Therefore, the access rule is

$$
P\land Q.
$$

Here,

$$
P=\text{valid password},
$$

and

$$
Q=\text{two-factor authentication completed}.
$$

\end{exampleblock}

\end{frame}

% ---------------------------------------------------------
% ANSWER 1 -- PART 2
% ---------------------------------------------------------

\begin{frame}{Propositions -- Question 1: Answer (Continued)}

\begin{exampleblock}{Step 2: Evaluate the Rule}

The given values are

$$
P=\text{True},
\qquad
Q=\text{False}.
$$

Therefore,

$$
P\land Q
=
\text{True}\land\text{False}
=
\text{False}.
$$

\end{exampleblock}

\end{frame}

% ---------------------------------------------------------
% ANSWER 1 -- PART 3
% ---------------------------------------------------------

\begin{frame}{Propositions -- Question 1: Answer (Conclusion)}

\begin{exampleblock}{Conclusion}

Since the access condition is false, the user is not granted access.

$$
\boxed{\text{Access Denied}}
$$

\end{exampleblock}

\end{frame}

% =========================================================
% QUESTION 2
% =========================================================

\begin{frame}{Propositions -- Question 2}

\begin{block}{Question}

An intrusion detection system raises an alert when either a suspicious
login location or an unusual login time is detected.

Let

$$
P:\text{suspicious login location is detected},
$$

$$
Q:\text{unusual login time is detected}.
$$

The system records that the login occurred from a suspicious location,
but at a normal time.

\end{block}

\end{frame}

\begin{frame}{Propositions -- Question 2 (Continued)}

\begin{block}{Question}

\begin{enumerate}
\item Express the alert condition as a proposition.
\item Determine whether an alert is raised.
\end{enumerate}

\end{block}

\end{frame}

% ---------------------------------------------------------
% ANSWER 2 -- PART 1
% ---------------------------------------------------------

\begin{frame}{Propositions -- Question 2: Answer}

\begin{exampleblock}{Step 1: Represent the Rule}

The system raises an alert when either condition occurs.

Therefore, the alert condition is

$$
P\lor Q.
$$

\end{exampleblock}

\end{frame}

% ---------------------------------------------------------
% ANSWER 2 -- PART 2
% ---------------------------------------------------------

\begin{frame}{Propositions -- Question 2: Answer (Continued)}

\begin{exampleblock}{Step 2: Evaluate the Conditions}

The recorded values are

$$
P=\text{True},
\qquad
Q=\text{False}.
$$

Thus,

$$
P\lor Q
=
\text{True}\lor\text{False}
=
\text{True}.
$$

\end{exampleblock}

\end{frame}

% ---------------------------------------------------------
% ANSWER 2 -- PART 3
% ---------------------------------------------------------

\begin{frame}{Propositions -- Question 2: Answer (Conclusion)}

\begin{exampleblock}{Conclusion}

Since the alert condition is true, the intrusion detection system
raises an alert.

$$
\boxed{\text{Alert Raised}}
$$

\end{exampleblock}

\end{frame}

% =========================================================
% QUESTION 3
% =========================================================

\begin{frame}{Propositions -- Question 3}

\begin{block}{Question}

A machine-learning service accepts a model for deployment only if the
model has completed testing and has an accuracy of at least $90%$.

Let

$$
P:\text{the model has completed testing},
$$

$$
Q:\text{the model has accuracy at least }90\%.
$$

A model has completed testing, but its accuracy is only $87%$.

\end{block}

\end{frame}

\begin{frame}{Propositions -- Question 3 (Continued)}

\begin{block}{Question}

Determine whether the model can be deployed and justify your answer
using a compound proposition.

\end{block}

\end{frame}

% ---------------------------------------------------------
% ANSWER 3 -- PART 1
% ---------------------------------------------------------

\begin{frame}{Propositions -- Question 3: Answer}

\begin{exampleblock}{Step 1: Identify the Conditions}

Both conditions are required for deployment.

Therefore, the deployment condition is

$$
P\land Q.
$$

For the given model,

$$
P=\text{True}.
$$

\end{exampleblock}

\end{frame}

% ---------------------------------------------------------
% ANSWER 3 -- PART 2
% ---------------------------------------------------------

\begin{frame}{Propositions -- Question 3: Answer (Continued)}

\begin{exampleblock}{Step 2: Evaluate Accuracy}

The model has an accuracy of $87%$.

Since

$$
87\% < 90\%,
$$

the accuracy condition is false:

$$
Q=\text{False}.
$$

\end{exampleblock}

\end{frame}

% ---------------------------------------------------------
% ANSWER 3 -- PART 3
% ---------------------------------------------------------

\begin{frame}{Propositions -- Question 3: Answer (Continued)}

\begin{exampleblock}{Step 3: Evaluate the Proposition}

Using

$$
P=\text{True},
\qquad
Q=\text{False},
$$

we obtain

$$
P\land Q
=
\text{True}\land\text{False}
=
\text{False}.
$$

Therefore, the deployment condition is false.

\end{exampleblock}

\end{frame}

% ---------------------------------------------------------
% ANSWER 3 -- PART 4
% ---------------------------------------------------------

\begin{frame}{Propositions -- Question 3: Answer (Conclusion)}

\begin{exampleblock}{Conclusion}

The model does not satisfy both conditions required for deployment.

Hence,

$$
\boxed{\text{The model cannot be deployed.}}
$$

The compound proposition

$$
P\land Q
$$

represents a rule in which all required conditions must be satisfied.

\end{exampleblock}

\end{frame}

% =========================================================
% QUESTION 4
% =========================================================

\begin{frame}{Propositions -- Question 4}

\begin{block}{Question}

A cloud server is considered operational when it is either in the
\textit{Running} state or in the \textit{Standby} state.

Let

$$
P:\text{server is in the Running state},
$$

$$
Q:\text{server is in the Standby state}.
$$

A monitoring system reports that the server is neither running nor in
standby mode.

\end{block}

\end{frame}

\begin{frame}{Propositions -- Question 4 (Continued)}

\begin{block}{Question}

\begin{enumerate}
\item Represent the operational condition as a proposition.
\item Determine whether the server is operational.
\item State the negation of the operational condition.
\end{enumerate}

\end{block}

\end{frame}

% ---------------------------------------------------------
% ANSWER 4 -- PART 1
% ---------------------------------------------------------

\begin{frame}{Propositions -- Question 4: Answer}

\begin{exampleblock}{Step 1: Represent the Condition}

The server is operational when it is in either state.

Therefore,

$$
P\lor Q.
$$

The monitoring system reports

$$
P=\text{False},
\qquad
Q=\text{False}.
$$

\end{exampleblock}

\end{frame}

% ---------------------------------------------------------
% ANSWER 4 -- PART 2
% ---------------------------------------------------------

\begin{frame}{Propositions -- Question 4: Answer (Continued)}

\begin{exampleblock}{Step 2: Determine the State}

Using the given values,

$$
P\lor Q
=
\text{False}\lor\text{False}
=
\text{False}.
$$

Therefore, the server does not satisfy the operational condition.

\end{exampleblock}

\end{frame}

% ---------------------------------------------------------
% ANSWER 4 -- PART 3
% ---------------------------------------------------------

\begin{frame}{Propositions -- Question 4: Answer (Continued)}

\begin{exampleblock}{Step 3: Negate the Condition}

The operational condition is

$$
P\lor Q.
$$

Its negation is

$$
\neg(P\lor Q).
$$

Using De Morgan's law,

$$
\neg(P\lor Q)
\equiv
\neg P\land\neg Q.
$$

\end{exampleblock}

\end{frame}

% ---------------------------------------------------------
% ANSWER 4 -- PART 4
% ---------------------------------------------------------

\begin{frame}{Propositions -- Question 4: Answer (Conclusion)}

\begin{exampleblock}{Interpretation}

The expression

$$
\neg P\land\neg Q
$$

means that the server is neither in the Running state nor in the
Standby state.

Therefore,

$$
\boxed{\neg(P\lor Q)\equiv\neg P\land\neg Q}
$$

is the required negation.

\end{exampleblock}

\end{frame}

% =========================================================
% QUESTION 5
% =========================================================

\begin{frame}{Propositions -- Question 5}

\begin{block}{Question}

A data-processing pipeline performs a backup when the following rule is
satisfied:

$$
\text{Backup occurs if the dataset is modified and the storage system
is available.}
$$

Let

$$
P:\text{dataset is modified},
$$

$$
Q:\text{storage system is available},
$$

$$
R:\text{backup occurs}.
$$

\end{block}

\end{frame}

\begin{frame}{Propositions -- Question 5 (Continued)}

\begin{block}{Question}

The system records that the dataset was modified and the storage system
was unavailable.

\medskip

\begin{enumerate}
\item Express the backup rule as a proposition.
\item Determine whether the backup should occur.
\end{enumerate}

\end{block}

\end{frame}

% ---------------------------------------------------------
% ANSWER 5 -- PART 1
% ---------------------------------------------------------

\begin{frame}{Propositions -- Question 5: Answer}

\begin{exampleblock}{Step 1: Identify the Conditions}

The backup rule contains three propositions:

$$
P:\text{dataset is modified},
$$

$$
Q:\text{storage system is available},
$$

$$
R:\text{backup occurs}.
$$

\end{exampleblock}

\end{frame}

% ---------------------------------------------------------
% ANSWER 5 -- PART 2
% ---------------------------------------------------------

\begin{frame}{Propositions -- Question 5: Answer (Continued)}

\begin{exampleblock}{Step 2: Represent the Rule}

The dataset must be modified and the storage system must be available.

Therefore, the required condition is

$$
P\land Q.
$$

If this condition is satisfied, a backup occurs.

Thus,

$$
\boxed{(P\land Q)\rightarrow R}.
$$

\end{exampleblock}

\end{frame}

% ---------------------------------------------------------
% ANSWER 5 -- PART 3
% ---------------------------------------------------------

\begin{frame}{Propositions -- Question 5: Answer (Continued)}

\begin{exampleblock}{Step 3: Use the Recorded Values}

The dataset was modified.

Therefore,

$$
P=\text{True}.
$$

The storage system was unavailable.

Therefore,

$$
Q=\text{False}.
$$

\end{exampleblock}

\end{frame}

% ---------------------------------------------------------
% ANSWER 5 -- PART 4
% ---------------------------------------------------------

\begin{frame}{Propositions -- Question 5: Answer (Continued)}

\begin{exampleblock}{Step 4: Evaluate the Antecedent}

The antecedent is

$$
P\land Q.
$$

Substituting the recorded values,

$$
P\land Q
=
\text{True}\land\text{False}
=
\text{False}.
$$

Therefore,

$$
\boxed{P\land Q=\text{False}}.
$$

\end{exampleblock}

\end{frame}

% ---------------------------------------------------------
% ANSWER 5 -- PART 5
% ---------------------------------------------------------

\begin{frame}{Propositions -- Question 5: Answer (Continued)}

\begin{exampleblock}{Step 5: Evaluate the Implication}

The complete proposition is

$$
(P\land Q)\rightarrow R.
$$

Since

$$
P\land Q=\text{False},
$$

the proposition becomes

$$
\text{False}\rightarrow R.
$$

An implication with a false antecedent is true.

Therefore,

$$
\boxed{(P\land Q)\rightarrow R=\text{True}}.
$$

\end{exampleblock}

\end{frame}

% ---------------------------------------------------------
% ANSWER 5 -- PART 6
% ---------------------------------------------------------

\begin{frame}{Propositions -- Question 5: Answer (Conclusion)}

\begin{exampleblock}{Interpretation}

The truth of the implication does not mean that a backup actually
occurred.

The required condition

$$
P\land Q
$$

was false because the storage system was unavailable.

Therefore, the backup rule was not violated.

\end{exampleblock}

\end{frame}

\begin{frame}{Propositions -- Question 5: Answer (Conclusion)}

\begin{exampleblock}{Final Result}

The system cannot require a backup when the storage system is
unavailable.

Hence,

$$
\boxed{\text{The backup rule is satisfied, but no backup is required.}}
$$

\end{exampleblock}

\end{frame}
